% !TEX root = ../main.tex
% Threats to Validity
% Claims to support: single dataset, prompt sensitivity, deterministic splits

\section{Threats to Validity}
\label{sec:threats}

\myitparagraph{Internal Validity.} 
Empirical selection of \bragtag's margin can potentially be a threat to internal validity, since the value could be specifically tuned for the models and test data in this study. To mitigate this, margin calibration used only the training split and did not involve the test split (\Cref{sec:ragtag}). Furthermore, the same fixed margin is used across all models and projects with no per-project or per-model tuning. Finally, the selected margin in this study is not framed as a universally optimal constant. Any project can re-derive the margin with negligible cost by only using the retrieval index of available historical issue reports.
The choice of hyperparameters and prompt format can also pose an internal threat to validity. Different hyperparameters and prompts can change model performance. We address this when reproducing the fine-tuning baseline in this study by adopting the same hyperparameters and prompt formats used in prior work~\cite{heo2025study,aracena2025applying}. 



\myitparagraph{External Validity.}
One threat to external validity concerns the generalizability of the results. Issue report text can vary substantially in writing style, vocabulary, and structure across software projects~\cite{akhavan2026linkanchorautonomousllmbasedagent}, which may threaten the generalizability of our findings. To mitigate this threat, we used a large benchmark of 6,600 issues from 11 different projects, each with its own reporting conventions and writing styles.
Another threat to external validity is the generalizability across settings. The primary aim of this work is to empirically assess the effectiveness of retrieval-augmented few-shot prompting versus LoRA fine-tuning for IRC and to explore how results vary with different model scales. This goal required exploring a large variable space across model sizes, $k$ values, data scopes, and different retrieval-based methods. However, Many configurations, including LLM architecture, prompt structure, embedding model, and retrieval technique, can each lead to different results, which can be a potential threat to the generalizability of our findings beyond the settings used in this study. Recognizing this threat, we encourage future work to build on our findings and conduct dedicated studies under different settings. To facilitate such studies, we designed our experiment pipeline in a modular way that allows all configurations to be easily substituted.

\myitparagraph{Construct Validity.}
The inherent randomness of LLM outputs can pose a potential threat to construct validity. The reported performance metrics can vary between different runs due to LLM randomness. To mitigate this variance, we ran the experiments three times and report the average.

\myparagraph{Reliability.} \fixme{Provide the link to replication package here.} 